\section{稳健性与敏感性分析}

\subsection{预测、识别与传导探针}

问题一将六类预测器置于78个相同目标月比较；no-change在三个期限均取得最低RMSE，Theta虽接近，
DM--HLN检验未显示稳定优势。事件CAR改用40个前期波动率、价格水平和GPR状态最接近的安慰剂，
最小可实现经验$p$值为$1/41=0.0244$，避免普通周末样本造成虚假精度。结构VAR的6、12、18、
24阶候选均检查稳定性；6阶模型白噪声检验通过，替代排序下关键价格冲击相关系数不低于0.999。

问题二使用联合最大$t$置信带、逐点区间、FDR校正、ARDL、疫情剔除和冲击正负拆分。多数PPI、
CPI和IAV联合区间跨零，GDP有效观测不足，因此稳健性分析的作用是限制结论强度，而不是事后挑选
显著规格。问题三的四个年度缓冲变量因未经审计统一判为“主分析失败”，不以替代规格恢复定量结论。

\subsection{尾部风险回测}

FHS与高斯随机游走使用相同的47个月度原点和1、3、6个月期限。除分位损失和覆盖率外，补充80\%、
90\%区间分数与近似CRPS\cite{gneiting2007}。长度6的配对块bootstrap显示：FHS仅在3个月两档区间分数和6个月80\%
区间分数上显著更优，其余差异为结论不充分。由此将“FHS全面优于基线”列为禁止表述。

\subsection{SAPR规则敏感性}

\begin{table}[htbp]
  \centering
  \small
  \caption{SAPR膝点对阈值、权重和块长度的敏感性}
  \label{tab:q4-sensitivity}
  \begin{tabular}{lcccc}
    \toprule
    变体 & $(\rho_N,\rho_S,\rho_E)$ & Pareto数 & 相对默认变化 & 检验期非支配 \\
    \midrule
    默认75/95阈值 & $(0.30,0.10,0.10)$ & 119 & 否 & 是 \\
    70/90阈值 & $(0.30,0.15,0.10)$ & 82 & 是 & 是 \\
    80/97.5阈值 & $(0.30,0.10,0.10)$ & 209 & 否 & 是 \\
    CPI权重增加20\% & $(0.30,0.10,0.10)$ & 121 & 否 & 是 \\
    IAV权重增加20\% & $(0.30,0.10,0.10)$ & 116 & 否 & 是 \\
    bootstrap块3/6/12 & $(0.30,0.10,0.10)$ & 119 & 否 & 是 \\
    \bottomrule
  \end{tabular}
\end{table}

除较低状态阈值外，膝点规则对权重、阈值上界和块长度保持不变；70/90阈值把压力期传导率提高到
0.15，说明普通与压力状态的划分是最敏感设计项。所有变体在检验期点估计上保持非支配，但正文的
“支持”判断还依赖配对目标差和硬约束，而不只依赖非支配概率。
